\scanpage{227}
Quelques calculs d'homotopie

Soit $k$ un anneau de base (comm). $\Phi_*=\Phi_*^k$ \deleted{\unsure{$k$}} le $k$-module $\tfrac12$ simplicial $I\longmapsto k^I$ $k_*$ le $k$-module $\tfrac12$ \unsure{simpl|simpl.} const défini par $k$, $I\longmapsto k$. On a une application diagonale
\[
 k_*\longrightarrow\Phi_*
\]
d'où un quotient
\[
 \Psi_*=\Phi_*/k_*\colon I\longmapsto k^I/k
\]

\uline{Proposition} 1 \quad $\Phi_*$ est homotope à 0 \deleted{\unsure{$0_k$}} (\inserted{i.e. au $k$-}\deleted{module} $\tfrac12$ simplicial 0)

Par le \unsure{topos|topo|yoga} de Puppe-Dold, il suffit de voir que le complexe de chaînes $\Phi_*^!$ associé \deleted{$\Phi$} à $\Phi$ est homo à 0. Pour ceci, on se rappelle que a
\[
 \Phi_n^{\unsure{!|\bullet}}=\{x\in\Phi_n\mid\partial_1^*(x)=\cdots=\partial_n^*(x)=0\}
\]
\deleted{d'où} aussitôt
\[
 \Phi_n^!=\begin{cases}0&\text{si }n\geq2\\ k&\text{si }n=0,1\end{cases}
\]
l'op. \deleted{\unsure{cob|co}} bord
\[
 \partial=\partial_0\colon\underset{\substack{\rotatebox{90}{$\simeq$}\\k}}{\Phi_1^!}\longrightarrow
 \underset{\substack{\rotatebox{90}{$\simeq$}\\k}}{\Phi_0^!}
\]
étant l'identité, d'où la conclusion

\uline{Corollaire 1.1} Pour tt foncteur $F\colon\operatorname{Mod}(k)\longrightarrow C$ ($C$ une catégorie) $F(\Phi_*)$ est homotope à $F(0)$. En particulier
\[
 \left\{\begin{aligned}
 \Lambda_k^{\unsure{*|i}}\Phi_*\,\inserted{(\Gamma_k^i\Phi_*,\operatorname{Sym}_k^i\Phi_*)}
 &\text{ est homotope à }\Lambda_k^i(0)\,\inserted{(\Gamma_k^i(0),\operatorname{Sym}_k^i(0))}
 =\begin{cases}0&\text{si }i\ne0\\(k)_*&\text{si }i=0\end{cases}\\
 \deleted{\Gamma_k^i\Phi_*}&\deleted{\longrightarrow\Gamma_k^i(0)=
 \left\{\substack{0\ \text{si }i\ne0\\k_*\ \text{si }i=0}\right.}
 \end{aligned}\right.
\]

\uline{Corollaire} 1.2. Les groupes d'homotopie de $\Psi_*=\Phi_*/k_*$ sont
\[
 \pi_i(\Psi_*\,\deleted{/k_*})\simeq\pi_{i-1}(k_*)\simeq
 \begin{cases}0&\text{si }i\ne1\\ k&\text{si }i=1\end{cases}
\]
En d'autres termes, $\Psi_*\deleted{/k_*}$ est un $K_*(k,1)$
\margin{$\big[$}
[En fait, on a $(\Psi_*^!\deleted{/k_*^!})^!=$ complexe réduit à $k$ placé en degré 1.

On conclut de ceci:\deleted{\unsure{et du topos de décalage que|et du yoga de Dold-Puppe}}

\uline{Corollaire} 1.\unsure{3|2}. \quad $\pi_j(\Lambda^i(\Phi_*/k_*))\simeq\begin{cases}0&\text{si }j\ne i\\k&\text{si }j=i\end{cases}$

(i.e. $\Lambda^i\Phi_*/k_*$ est un $K_*(k,i)$).

% archive sheet number: 226

